\documentclass{article}%
\usepackage[T1]{fontenc}%
\usepackage[utf8]{inputenc}%
\usepackage{lmodern}%
\usepackage{textcomp}%
\usepackage{lastpage}%
\usepackage{geometry}%
\geometry{margin=2cm}%
\usepackage{expex}%
\usepackage[T1]{fontenc}%
\usepackage[utf8]{inputenc}%
\usepackage[spanish]{babel}%
%
\title{Análisis Lingüístico: Gramatica Normativa Kaqchikel}%
\author{Generado por el TFM de Elías Carrasco}%
%
\begin{document}%
\normalsize%
\maketitle%
\section{Page 1}%
\label{sec:Page1}%
INC{-}B3 %
A3{-}querer{-}SC %
Lo quiere %
\subsection{Negación de una acción}%
\label{subsec:Negacindeunaaccin}%

%
\par\medskip%
\ex
\textit{No lo quiere} \\
\begingl
\gla k'am //
\glb k'am chi-ø y-oche-j //
\glc \textbf{NEG} \textbf{INC-}\textbf{B3} \textbf{A3-}\textbf{querer-}\textbf{SC} //
\glft \textbf{Sintaxis:} [k'am]\textsubscript{SN} //
\endgl
\xe
%
\medskip%
\subsection{Oración base (acción en potencial)}%
\label{subsec:Oracinbase(accinenpotencial)}%

%
\par\medskip%
\ex
\textit{Él se bañará} \\
\begingl
\gla hoq ø achn-oq //
\glb hoq ø achn-oq //
\glc \textbf{POT-} \textbf{B3-} \textbf{bañar-}\textbf{SC} //
\endgl
\xe
%
\medskip%
\subsection{Negación de una acción}%
\label{subsec:Negacindeunaaccin}%

%
\par\medskip%
\ex
\textit{No se bañará} \\
\begingl
\gla man hoq ø achn-oq naq //
\glb man hoq ø achn-oq naq //
\glc \textbf{no} \textbf{POT-B3} \textbf{bañar-SC} \textbf{PRO:él} \textbf{PRO:él} //
\endgl
\xe
%
\medskip%
\subsection{Oración base (acción en incompletivo)}%
\label{subsec:Oracinbase(accinenincompletivo)}%

%
\par\medskip%
\ex
\textit{Me iré mañana} \\
\begingl
\gla ch-in toj ir //
\glb ch-in toj ir //
\glc \textbf{INC-}\textbf{B1SG} \textbf{ir} \textbf{mañana} //
\glft \textbf{Sintaxis:} [ch-in]\textsubscript{SN} //
\endgl
\xe
%
\medskip%
\subsection{Negación de una acción}%
\label{subsec:Negacindeunaaccin}%

%
\par\medskip%
\ex
\textit{No me iré mañana} \\
\begingl
\gla maj hin toq ir //
\glb maj hin toq ir //
\glc \textbf{no} \textbf{A1SG} \textbf{ir} \textbf{mañana} //
\endgl
\xe
%
\medskip%
\subsection{Oración base (predicado no verbal)}%
\label{subsec:Oracinbase(predicadonoverbal)}%

%
\par\medskip%
\ex
\textit{Ella está parada} \\
\begingl
\gla lek-an ø aj ix //
\glb lek-an ø aj ix //
\glc \textbf{parado-POS} \textbf{B3} \textbf{DIR:hacia arriba} \textbf{PRO:ella} //
\glft \textbf{Sintaxis:} [lek-an]\textsubscript{SN} //
\endgl
\xe
%
\medskip%
\subsection{Negación de un predicado no verbal}%
\label{subsec:Negacindeunpredicadonoverbal}%

%
\par\medskip%
\ex
\textit{} \\
\begingl
\gla man lekan-oq ø aj ix //
\glb man lekan-oq ø aj ix //
\glc \textbf{no} \textbf{} \textbf{} \textbf{} \textbf{} //
\endgl
\xe
%
\medskip

%
\end{document}